\documentclass[12pt, a4paper]{article}

% --- PACOTES ESSENCIAIS ---
\usepackage[utf8]{inputenc}
\usepackage[brazil]{babel}
\usepackage[top=3cm, bottom=2cm, left=3cm, right=2cm]{geometry}
\usepackage[final]{graphicx}
\usepackage{xcolor}
\usepackage{microtype}
\usepackage{tcolorbox}
\usepackage{enumitem}
\usepackage{caption}
\usepackage{hyperref}

\hypersetup{
    colorlinks=true,
    linkcolor=black,
    citecolor=black,
    urlcolor=blue
}

\pagestyle{plain}

\begin{document}

% ==============================================================================
% CAPA
% ==============================================================================
\begin{titlepage}
    \begin{center}
        \includegraphics[height=2.2cm, keepaspectratio]{Logo_UTFPR.png}\\[1cm]
        \textbf{\large UNIVERSIDADE TECNOLÓGICA FEDERAL DO PARANÁ}\\
        \textbf{\large CAMPUS CORNÉLIO PROCÓPIO}\\[0.3cm]
        \hrule height 1pt
        \vspace{3.5cm}
        {\Huge \bfseries RELATÓRIO}\\[0.5cm]
        {\Large \bfseries 4ª OFICINA DE ORIGAMI SEISHUN}\\[4.5cm]
        \vfill
        \textbf{\large CORNÉLIO PROCÓPIO}\\
        \textbf{\large 2026}
    \end{center}
\end{titlepage}

% ==============================================================================
% SUMÁRIO
% ==============================================================================
\tableofcontents
\newpage

% ==============================================================================
% CONTEÚDO DO RELATÓRIO
% ==============================================================================

\section{Introdução}
O presente relatório visa documentar as atividades e resultados obtidos durante a 4ª Oficina de Origami Seishun. A ação buscou proporcionar aos participantes uma introdução prática e teórica à arte secular do origami, promovendo o desenvolvimento de habilidades motoras finas, capacidade de concentração, espacialidade e criatividade através da confecção de figuras em papel.

A oficina foi estruturada metodologicamente em três ``Ilhas'' temáticas, cada uma correspondente a um nível de dificuldade progressivo (Fácil, Médio e Difícil). Esta divisão permitiu atender de forma inclusiva tanto os participantes iniciantes quanto aqueles com experiência prévia na arte das dobraduras.

\section{Métodos e Procedimentos}

\begin{itemize}[leftmargin=*]
    \item \textbf{Ilhas de Origami:} Organização do espaço físico em três estações distintas baseadas no nível de complexidade das dobraduras (Fácil, Médio e Difícil), possibilitando a livre escolha do participante de acordo com sua aptidão e interesse.
    \item \textbf{Monitores Dedicados:} Atuação de 1 a 3 monitores por ilha, garantindo atendimento individualizado, mediação de dúvidas técnicas e auxílio contínuo no acompanhamento dos diagramas.
    \item \textbf{Passo a Passo Visual:} Disponibilização de diagramas e guias visuais instrucionais em todas as bancadas para orientação autônoma dos participantes.
    \item \textbf{Material para Treino:} Fornecimento contínuo de papel sulfite padronizado para experimentação e treino das técnicas de dobra antes da execução final.
    \item \textbf{Confecção e Entrega da Peça Final:} Disponibilização de papéis especiais para origami ao término da prática para que cada participante pudesse produzir e levar sua peça final para casa.
\end{itemize}

\begin{figure}[htbp]
    \centering
    \includegraphics[height=7cm, keepaspectratio]{1.png}
    \caption{Visão geral da sala durante a realização da oficina.}
    \label{fig:foto1}
\end{figure}

\begin{figure}[htbp]
    \centering
    \includegraphics[height=7cm, keepaspectratio]{2.png}
    \caption{Exemplos de origamis confeccionados pelos participantes.}
    \label{fig:foto2}
\end{figure}

\subsection{Equipe Executora}
A equipe responsável pela organização e execução das atividades foi composta por docentes e discentes voluntários, conforme detalhado a seguir:

\begin{itemize}[leftmargin=*]
    \item Prof. Dr. André Luís Shiguemoto -- Coordenador do Projeto;
    \item Profa. Dra. Gabriela Helena Bauab Shiguemoto -- Vice-Coordenadora do Projeto;
    \item André Luiz da Silva Junior, Voluntário;
    \item Ayumi Wassano, Voluntária;
    \item Bianca Natsumi Nishimura, Voluntária;
    \item João Vitor Vieira Ribeiro Doneda, Voluntário;
    \item Thalita Yukare Pereira Kinoshita, Voluntária.
\end{itemize}

\subsection*{Equipe de Apoio}
\begin{itemize}[leftmargin=*]
    \item Anna Liz Michlovská R. Rodrigues, Voluntária;
    \item Gustavo Bashio Nakano, Voluntário;
\end{itemize}

\section{Resultados e Discussões}
A 4ª edição do evento registrou um aumento expressivo no número de participantes em relação às edições anteriores. Destaca-se que a grande maioria dos presentes relatou estar participando da atividade pela primeira vez. 

Observou-se uma preferência acentuada dos participantes pelas atividades dispostas na ilha de nível fácil. As figuras com maior adesão e interesse do público foram: o tradicional *Tsuru*, a estrela ninja (*Shuriken*), flores, tartarugas, baleias, sapos saltadores e corações. A Figura~\ref{fig:foto1} ilustra o engajamento do público na sala, enquanto a Figura~\ref{fig:foto2} apresenta alguns dos modelos produzidos durante a oficina.

\section{Conclusões}
A 4ª Oficina de Origami Seishun cumpriu satisfatoriamente seus objetivos, mantendo a consistência e o fluxo metodológico consolidado nas edições anteriores. O aumento do público e o interesse dos novos participantes ratificam a relevância da atividade como espaço de integração comunitária, aprendizado e difusão cultural no campus.

\end{document}
